% !PRINTED 340
\weberpage{340}
% running head: 340 | Elfter Abschnitt. | §. 91.
%<<CONT
according, then, as $T$ is positive or negative, the first or the
second value of $r$ turns out real, and in essentials everything
remains as before.

From this results the theorem:

\gesperrt{Every real ternary cubic form with non-vanishing
discriminant can be brought by a real
transformation to the real canonical form.}

One can now easily give the coordinates of the points of inflection
in the canonical coordinate system. One obtains them
by setting, in the equation
\[
y_1^3+y_2^3+y_3^3+6m\,y_1y_2y_3=0
\]
$y_1$, $y_2$, $y_3=0$ in turn. Since one of the non-vanishing
coordinates can be set $=1$, one finds, if $\varepsilon$
is an imaginary third root of unity, the following values
of the coordinates $(y_1, y_2, y_3)$:
\weq{16}{\begin{array}{ccc}
(0, 1, -1), & (0, 1, -\varepsilon), & (0, 1, -\varepsilon^2), \\
(-1, 0, 1), & (-\varepsilon, 0, 1), & (-\varepsilon^2, 0, 1), \\
(1, -1, 0), & (1, -\varepsilon, 0), & (1, -\varepsilon^2, 0),
\end{array}}
and it follows from this that \gesperrt{always three of the inflexion points
are real and six imaginary.}

In the table (16), the three points standing in any one row
lie on a straight line, namely on the lines
\[
y_1=0, \quad y_2=0, \quad y_3=0.
\]

It is easy, after this, to assemble also the remaining triples of points
that lie on a straight line, and thereby
to determine these twelve lines. One has only always to
assemble such points whose coordinates give a vanishing
determinant. One can arrange these twelve lines
in triples in four ways, so that in each triple all
nine inflexion points occur. The first of these triples is
formed by the rows of (16). The three others read thus:
\weq{17}{\begin{array}{ccc}
(0, 1, -1), & (-1, 0, 1), & (1, -1, 0), \\
(0, 1, -\varepsilon), & (-\varepsilon, 0, 1), & (1, -\varepsilon, 0), \\
(0, 1, -\varepsilon^2), & (-\varepsilon^2, 0, 1), & (1, -\varepsilon^2, 0),
\end{array}}
%CONT>>
